\documentclass[11pt, a4paper]{article}

% bfw-style.tex — gemeinsame Style-Definitionen für alle Handouts/Aufgabenhefte
% Voraussetzung: Kompilation mit xelatex (wegen fontspec)

\usepackage[ngerman]{babel}
\usepackage{geometry}
\geometry{a4paper, top=2.2cm, bottom=2.2cm, left=2.3cm, right=2.3cm}

\usepackage{fontspec}
% Fallback-Kette: Source Sans Pro -> Calibri -> Segoe UI -> Latin Modern Sans
% (Latin Modern ist immer mit MiKTeX vorhanden)
\IfFontExistsTF{Source Sans Pro}{%
  \setmainfont{Source Sans Pro}[Ligatures=TeX]%
}{\IfFontExistsTF{Calibri}{%
  \setmainfont{Calibri}[Ligatures=TeX]%
}{\IfFontExistsTF{Segoe UI}{%
  \setmainfont{Segoe UI}[Ligatures=TeX]%
}{%
  \setmainfont{Latin Modern Sans}[Ligatures=TeX]%
}}}
\IfFontExistsTF{Source Code Pro}{%
  \setmonofont{Source Code Pro}[Scale=0.92]%
}{\IfFontExistsTF{Consolas}{%
  \setmonofont{Consolas}[Scale=0.92]%
}{%
  \setmonofont{Latin Modern Mono}[Scale=0.92]%
}}

\usepackage{xcolor}
% Farbpalette: hell, freundlich, modern
\definecolor{bfwPrimary}{HTML}{0EA5A4}    % Petrol/Türkis - Primär
\definecolor{bfwPrimaryDark}{HTML}{0F766E} % dunkler Petrol für Headings
\definecolor{bfwAccent}{HTML}{F59E0B}     % Amber - Akzent
\definecolor{bfwSuccess}{HTML}{10B981}    % Grün - Erfolg / Lösung
\definecolor{bfwWarn}{HTML}{F97316}       % Orange - Achtung
\definecolor{bfwInk}{HTML}{0F172A}        % fast Schwarz
\definecolor{bfwInkSoft}{HTML}{475569}    % gedämpftes Grau
\definecolor{bfwBgSoft}{HTML}{F8FAFC}     % off-white
\definecolor{bfwBgTeal}{HTML}{ECFEFF}     % helles Türkis
\definecolor{bfwBgAmber}{HTML}{FEF3C7}    % helles Gelb
\definecolor{bfwBgRose}{HTML}{FFE4E6}     % helles Rosé für Eselsbrücken

\usepackage[colorlinks=true, linkcolor=bfwPrimaryDark, urlcolor=bfwPrimaryDark]{hyperref}
\usepackage{enumitem}
\setlist[itemize]{leftmargin=*, itemsep=2pt, topsep=2pt}
\setlist[enumerate]{leftmargin=*, itemsep=2pt, topsep=2pt}

\usepackage[most]{tcolorbox}
\usepackage{booktabs}
\usepackage{tikz}
\usetikzlibrary{decorations.pathmorphing, positioning, shapes.geometric, arrows.meta}

% Merkkästchen — wichtige Definition / Kernpunkt
\newtcolorbox{merksatz}[1][]{%
  enhanced, breakable,
  colback=bfwBgTeal,
  colframe=bfwPrimary,
  boxrule=0.6pt,
  arc=4pt,
  left=10pt, right=10pt, top=8pt, bottom=8pt,
  fonttitle=\bfseries\color{bfwPrimaryDark},
  title={\faLightbulb\ Merksatz},
  #1
}

% Eselsbrücke — Mnemoniks
\newtcolorbox{eselsbruecke}[1][]{%
  enhanced, breakable,
  colback=bfwBgRose,
  colframe=bfwAccent,
  boxrule=0.6pt,
  arc=4pt,
  left=10pt, right=10pt, top=8pt, bottom=8pt,
  fonttitle=\bfseries\color{bfwWarn},
  title={Eselsbrücke},
  #1
}

% Beispiel-Box
\newtcolorbox{beispiel}[1][]{%
  enhanced, breakable,
  colback=bfwBgSoft,
  colframe=bfwInkSoft,
  boxrule=0.4pt,
  arc=3pt,
  left=10pt, right=10pt, top=8pt, bottom=8pt,
  fonttitle=\bfseries\color{bfwInk},
  title={Beispiel},
  #1
}

% Lösungs-Box (für Aufgabenhefte)
\newtcolorbox{loesung}[1][]{%
  enhanced, breakable,
  colback=bfwBgAmber!40,
  colframe=bfwSuccess,
  boxrule=0.6pt,
  arc=3pt,
  left=10pt, right=10pt, top=8pt, bottom=8pt,
  fonttitle=\bfseries\color{bfwSuccess!70!black},
  title={Lösung},
  #1
}

% Glossar-Eintrag
\newcommand{\glossareintrag}[2]{%
  \par\noindent\textbf{\color{bfwPrimaryDark}#1} \enspace #2\par\smallskip
}

% Section-Stil — etwas farbiger als Standard
\usepackage{titlesec}
\titleformat{\section}
  {\Large\bfseries\color{bfwPrimaryDark}}
  {\thesection}{0.6em}{}
\titleformat{\subsection}
  {\large\bfseries\color{bfwPrimary}}
  {\thesubsection}{0.5em}{}
\titleformat{\subsubsection}
  {\normalsize\bfseries\color{bfwInk}}
  {\thesubsubsection}{0.4em}{}

% TikZ-Stil "skizzenhaft" — etwas weicher als Rough.js, aber stabil
\tikzset{
  roughbox/.style = {
    draw=bfwPrimaryDark, thick, fill=bfwBgTeal,
    rounded corners=4pt, inner sep=6pt
  },
  roughaccent/.style = {
    draw=bfwAccent, thick, fill=bfwBgAmber,
    rounded corners=4pt, inner sep=6pt
  },
  roughline/.style = {
    draw=bfwInkSoft, thick, -{Stealth[length=2.5mm]}
  }
}

% Bullet-Symbole (bewusst kein FontAwesome — vermeidet Font-Installations-Probleme)
\newcommand{\faLightbulb}{$\bullet$}
\newcommand{\faBookmark}{$\bullet$}

% Header/Footer
\usepackage{fancyhdr}
\pagestyle{fancy}
\fancyhf{}
\renewcommand{\headrulewidth}{0pt}
\fancyfoot[L]{\small\color{bfwInkSoft}\thepage}
\fancyfoot[R]{\small\color{bfwInkSoft} BFW M-064 Beschaffung im Büromanagement}


\title{\vspace*{-1.5cm}\Huge\bfseries\color{bfwPrimaryDark}Aufgabenheft Session~2\\[0.4em]
  \large\color{bfwInkSoft}\textnormal{Bedürfnisse, Bedarf \& Wirtschaftskreislauf --- Übungen im IHK-Klausurformat}}
\author{\textsf{Lernfeld 1 --- Die eigene Rolle im Betrieb mitgestalten \textperiodcentered{} \small\copyright\ Can Siebert\,\textperiodcentered\,2026}}
\date{Selbstlernmaterial}

\begin{document}
\maketitle
\thispagestyle{empty}

\begin{merksatz}[title={\bfseries So nutzen Sie dieses Aufgabenheft}]
Bearbeiten Sie alle Aufgaben zunächst \emph{ohne} in die Lösungen zu schauen. Schreiben
Sie Ihre Antworten in vollständigen Sätzen --- die IHK-Prüfung verlangt formulierte
Begründungen, nicht bloße Stichworte. Beachten Sie die \emph{Operatoren} (nennen,
erläutern, beurteilen, begründen) --- sie geben den geforderten Antwort-Tiefegrad vor.
Erst danach öffnen Sie den Lösungsteil ab Seite~\pageref{loesungen} und vergleichen.
\end{merksatz}

\tableofcontents
\vspace{1em}
\hrule

\section{Aufgaben}

\subsection*{Aufgabe~1 --- Bedürfnisarten (10 Punkte)}

\noindent\textbf{\color{bfwAccent}YouTube-Vorbereitung:}\enspace \href{https://www.youtube.com/results?search_query=Bed%C3%BCrfnisarten+Existenzbed%C3%BCrfnis+Kulturbed%C3%BCrfnis+Luxusbed%C3%BCrfnis}{Bedürfnisarten einfach erklärt}

\medskip
\begin{enumerate}[label=(\alph*)]
  \item \textbf{Erläutern} Sie den Unterschied zwischen Existenz-, Kultur- und Luxusbedürfnissen. \hfill (3~P.)
  \item \textbf{Ordnen} Sie zu (Individual- oder Kollektivbedürfnis): eigenes Smartphone, sauberes Trinkwasser, gutes Autobahnnetz. \hfill (3~P.)
  \item \textbf{Erklären} Sie den Unterschied zwischen einem akuten und einem latenten Bedürfnis. \hfill (2~P.)
  \item \textbf{Beurteilen} Sie: Warum sind die Grenzen zwischen den Bedürfnisarten oft fließend? \hfill (2~P.)
\end{enumerate}

\subsection*{Aufgabe~2 --- Bedürfnis, Bedarf, Nachfrage (12 Punkte)}

\noindent\textbf{\color{bfwAccent}YouTube-Vorbereitung:}\enspace \href{https://www.youtube.com/results?search_query=Bed%C3%BCrfnis+Bedarf+Nachfrage+einfach+erkl%C3%A4rt}{Bedürfnis --- Bedarf --- Nachfrage}

\medskip
\begin{enumerate}[label=(\alph*)]
  \item \textbf{Bringen} Sie in die richtige Reihenfolge: Nachfrage, Bedürfnis, Bedarf. \hfill (3~P.)
  \item \textbf{Erläutern} Sie die Rolle der Kaufkraft beim Übergang vom Bedürfnis zum Bedarf. \hfill (3~P.)
  \item \textbf{Beschreiben} Sie an einem Beispiel den Weg vom Mangelgefühl bis zum Kauf. \hfill (3~P.)
  \item \textbf{Erklären} Sie die Grundidee der Preisbildung auf dem Markt. \hfill (3~P.)
\end{enumerate}

\subsection*{Aufgabe~3 --- Güterarten (15 Punkte)}

\noindent\textbf{\color{bfwAccent}YouTube-Vorbereitung:}\enspace \href{https://www.youtube.com/results?search_query=G%C3%BCterarten+freie+knappe+Konsumg%C3%BCter+Produktionsg%C3%BCter}{Güterarten im Überblick}

\medskip
\begin{enumerate}[label=(\alph*)]
  \item \textbf{Unterscheiden} Sie freie und knappe Güter mit je einem Beispiel. \hfill (3~P.)
  \item \textbf{Ordnen} Sie ein: Luft, Schreibtisch, Reparaturleistung, Patent, dienstliches Fahrzeug, Holz als Rohstoff. \hfill (6~P.)
  \item \textbf{Erläutern} Sie den Unterschied zwischen Komplementär- und Substitutionsgütern. \hfill (3~P.)
  \item \textbf{Erklären} Sie den Unterschied zwischen Gebrauchs- und Verbrauchsgut. \hfill (3~P.)
\end{enumerate}

\subsection*{Aufgabe~4 --- Ökonomisches Prinzip (12 Punkte)}

\noindent\textbf{\color{bfwAccent}YouTube-Vorbereitung:}\enspace \href{https://www.youtube.com/results?search_query=%C3%96konomisches+Prinzip+Maximalprinzip+Minimalprinzip}{Maximal- \& Minimalprinzip}

\medskip
\begin{enumerate}[label=(\alph*)]
  \item \textbf{Erklären} Sie Maximal- und Minimalprinzip mit je einem Satz. \hfill (3~P.)
  \item \textbf{Ordnen} Sie zu: (1) Mit fester Menge Holz möglichst viele Schreibtische bauen. (2) Bestellte Menge Schreibtische mit möglichst wenig Holz produzieren. \hfill (3~P.)
  \item \textbf{Beurteilen} Sie: Lassen sich beide Prinzipien gleichzeitig anwenden? Begründen Sie. \hfill (3~P.)
  \item \textbf{Erläutern} Sie, warum das ökonomische Prinzip aus der Knappheit der Güter folgt. \hfill (3~P.)
\end{enumerate}

\subsection*{Aufgabe~5 --- Produktionsfaktoren (12 Punkte)}

\noindent\textbf{\color{bfwAccent}YouTube-Vorbereitung:}\enspace \href{https://www.youtube.com/results?search_query=Produktionsfaktoren+Boden+Arbeit+Kapital+Betriebsmittel+Werkstoffe}{Produktionsfaktoren erklärt}

\medskip
\begin{enumerate}[label=(\alph*)]
  \item \textbf{Nennen} Sie die drei volkswirtschaftlichen Produktionsfaktoren mit je einem Beispiel. \hfill (3~P.)
  \item \textbf{Unterscheiden} Sie Elementarfaktoren und dispositive Faktoren. \hfill (3~P.)
  \item \textbf{Ordnen} Sie zu (Roh-, Hilfs- oder Betriebsstoff): Holz, Leim, Strom. \hfill (3~P.)
  \item \textbf{Erklären} Sie, warum Kapital als derivativer (abgeleiteter) Faktor gilt. \hfill (3~P.)
\end{enumerate}

\subsection*{Aufgabe~6 --- Wirtschaftskreislauf (14 Punkte)}

\noindent\textbf{\color{bfwAccent}YouTube-Vorbereitung:}\enspace \href{https://www.youtube.com/results?search_query=Wirtschaftskreislauf+einfach+erweitert+vollst%C3%A4ndig+erkl%C3%A4rt}{Wirtschaftskreislauf Schritt für Schritt}

\medskip
\begin{enumerate}[label=(\alph*)]
  \item \textbf{Beschreiben} Sie Geld- und Güterstrom im einfachen Wirtschaftskreislauf. \hfill (4~P.)
  \item \textbf{Erläutern} Sie die Rolle des Sektors Banken im erweiterten Kreislauf. \hfill (3~P.)
  \item \textbf{Nennen} Sie vier Aufgaben des Staates im vollständigen Wirtschaftskreislauf. \hfill (4~P.)
  \item \textbf{Erklären} Sie, wodurch aus dem vollständigen ein offener Wirtschaftskreislauf wird. \hfill (3~P.)
\end{enumerate}

\vspace{1em}
\noindent\textbf{Punkte gesamt: 75}

\newpage

\section{Lösungen}\label{loesungen}

\subsection*{Lösung Aufgabe~1}
\begin{loesung}
\textbf{(a)} \emph{Existenzbedürfnisse} (Grundbedürfnisse) sichern das Überleben (Essen,
Trinken, Wohnung). \emph{Kulturbedürfnisse} werden vom Kulturkreis geprägt und als üblich
angesehen (Bildung, Kino, modische Kleidung). \emph{Luxusbedürfnisse} sind nicht notwendige
Ansprüche, die sich nicht jeder leisten kann (Yacht, Sportwagen, Villa).

\textbf{(b)} Eigenes Smartphone = \emph{Individualbedürfnis}; sauberes Trinkwasser und gutes
Autobahnnetz = \emph{Kollektivbedürfnisse} (entstehen aus dem Zusammenleben).

\textbf{(c)} Ein \emph{akutes} Bedürfnis ist ein konkreter, offener Wunsch. Ein
\emph{latentes} Bedürfnis schlummert und ist dem Menschen nicht bewusst; es kann z.\,B.\
durch Werbung geweckt werden.

\textbf{(d)} Weil dasselbe Gut je nach Person und Lebensumstand unterschiedlich eingeordnet
wird: Ein Smartphone kann für den einen Kulturbedürfnis, für den anderen Luxus sein. Eine
eindeutige Abgrenzung ist deshalb nicht immer möglich.
\end{loesung}

\subsection*{Lösung Aufgabe~2}
\begin{loesung}
\textbf{(a)} Richtige Reihenfolge: \textbf{Bedürfnis} $\rightarrow$ \textbf{Bedarf}
$\rightarrow$ \textbf{Nachfrage}.

\textbf{(b)} Erst wenn für ein Bedürfnis die finanziellen Mittel (Einkommen bzw.\ Kaufkraft)
vorhanden sind, wird aus dem Bedürfnis ein \emph{Bedarf}. Die Kaufkraft ist also das
Bindeglied; ohne Einkommen bleibt das Bedürfnis wirkungslos.

\textbf{(c)} Beispiel: Wunsch nach einem neuen Smartphone = Mangelgefühl (\emph{Bedürfnis})
$\rightarrow$ Ausbildungsvergütung/Einkommen macht den Kauf möglich (\emph{Bedarf})
$\rightarrow$ tatsächlicher Kauf im Onlineshop (\emph{Nachfrage}).

\textbf{(d)} Auf dem Markt treffen Angebot und Nachfrage zusammen. Ist ein Gut sehr knapp
oder die Nachfrage groß, \emph{steigt} der Preis; ist es reichlich vorhanden oder die
Nachfrage gering, \emph{sinkt} der Preis. Der Preis signalisiert so die Knappheit.
\end{loesung}

\subsection*{Lösung Aufgabe~3}
\begin{loesung}
\textbf{(a)} \emph{Freie Güter} sind in großer Menge vorhanden und haben keinen Preis
(z.\,B.\ Luft). \emph{Knappe Güter} (Wirtschaftsgüter) sind begrenzt, verursachen Kosten
und haben einen Preis (z.\,B.\ ein Schreibtisch).

\textbf{(b)} Luft = freies Gut; Schreibtisch = knappes, materielles Konsum-/Gebrauchsgut;
Reparaturleistung = Dienstleistung (immateriell); Patent = Recht (immateriell); dienstliches
Fahrzeug = Produktionsgut (Gebrauchsgut); Holz als Rohstoff = Produktions-/Verbrauchsgut.

\textbf{(c)} \emph{Komplementärgüter} ergänzen sich und werden gemeinsam nachgefragt
(Maschine \& Strom, Butter \& Toast); die Nachfrage entwickelt sich gleichläufig.
\emph{Substitutionsgüter} ersetzen sich (Butter/Margarine); die Nachfrage verläuft
gegenläufig.

\textbf{(d)} \emph{Gebrauchsgüter} können mehrmals verwendet werden (Maschine).
\emph{Verbrauchsgüter} nur einmal (Holz als Rohstoff in der Produktion).
\end{loesung}

\subsection*{Lösung Aufgabe~4}
\begin{loesung}
\textbf{(a)} \emph{Maximalprinzip:} Mit gegebenen Mitteln einen maximalen Erfolg erzielen.
\emph{Minimalprinzip:} Ein gegebenes Ziel mit minimalem Aufwand erreichen.

\textbf{(b)} (1) = \emph{Maximalprinzip} (gegebene Mittel = Holzmenge, maximaler Erfolg =
möglichst viele Schreibtische). (2) = \emph{Minimalprinzip} (gegebenes Ziel = bestellte
Menge, minimaler Aufwand = möglichst wenig Holz).

\textbf{(c)} Nein. Man kann immer nur eine Größe fest vorgeben und die andere optimieren:
entweder das Mittel (Maximalprinzip) oder das Ziel (Minimalprinzip). Beides gleichzeitig zu
maximieren und zu minimieren ist logisch nicht möglich.

\textbf{(d)} Weil die Bedürfnisse unbegrenzt, die Güter aber knapp sind und Kosten
verursachen. Aus diesem Spannungsverhältnis folgt, dass Haushalte und Unternehmen planvoll,
also nach dem ökonomischen Prinzip, mit den knappen Gütern umgehen müssen.
\end{loesung}

\subsection*{Lösung Aufgabe~5}
\begin{loesung}
\textbf{(a)} \emph{Boden} (z.\,B.\ Grundstück der Betriebsstätte), \emph{Arbeit} (geistige/
körperliche Tätigkeit gegen Einkommen), \emph{Kapital} (Geld-/Sachkapital, z.\,B.\ Maschinen).

\textbf{(b)} \emph{Elementarfaktoren}: ausführende Arbeitskräfte, Betriebsmittel (Maschinen,
Fahrzeuge, Werkzeuge) und Werkstoffe. \emph{Dispositive Faktoren}: Leitung, Planung,
Organisation und Überwachung (Kontrolle).

\textbf{(c)} Holz = \emph{Rohstoff} (Hauptbestandteil); Leim = \emph{Hilfsstoff}
(Nebenbestandteil); Strom = \emph{Betriebsstoff} (geht nicht in das Produkt ein).

\textbf{(d)} Kapital steht nicht von Anfang an zur Verfügung, sondern entsteht erst durch
Sparen (Konsumverzicht) und durch die Kombination der Faktoren Boden und Arbeit. Deshalb
gilt es als abgeleiteter (derivativer) Faktor.
\end{loesung}

\subsection*{Lösung Aufgabe~6}
\begin{loesung}
\textbf{(a)} \emph{Güterstrom}: Die Haushalte stellen den Unternehmen die
Produktionsfaktoren (Boden, Arbeit, Kapital) bereit; die Unternehmen liefern dafür
Konsumgüter und Dienstleistungen an die Haushalte zurück. \emph{Geldstrom}: Die Unternehmen
zahlen Einkommen/Entlohnung an die Haushalte; die Haushalte geben dieses Einkommen als
Konsumausgaben wieder an die Unternehmen. Beide Ströme laufen gegenläufig und gleichen sich
aus.

\textbf{(b)} Die \emph{Banken} nehmen die Ersparnisse der Haushalte auf und stellen sie den
Unternehmen als Kredite zur Verfügung. Damit können die Unternehmen investieren (z.\,B.\ in
neue Maschinen und Technologien). Im Modell entsprechen die Ersparnisse den Investitionen.

\textbf{(c)} Der \emph{Staat} bezieht Steuern von Haushalten und Unternehmen, zahlt Löhne/
Gehälter und Sozialleistungen, subventioniert Unternehmen und konsumiert selbst Güter und
Dienstleistungen.

\textbf{(d)} Wird zusätzlich das \emph{Ausland} einbezogen (Import und Export von Gütern,
Dienstleistungen und Zahlungen), spricht man vom Wirtschaftskreislauf einer offenen
Volkswirtschaft.
\end{loesung}

\end{document}
